\olchapter{pl}{syn}{Syntax and Semantics}

\begin{editorial}
  This is a very quick summary of definitions only. It should be
  expanded to provide a gentle intro to proofs by induction on
  formulas, with lots more examples.
\end{editorial}



\olfileid{pl}{syn}{int}

\olsection{Introduction}

Propositional logic deals with !!{formula}s that are built from
!!{propositional variable}s using the propositional connectives
$\lnot$, $\land$, $\lor$, $\lif$, and $\liff$.  Intuitively,
!!a{propositional variable}~$p$ stands for a sentence or proposition
that is true or false. Whenever the ``truth value'' of the
!!{propositional variable} in !!a{formula} is determined, so is the
truth value of any !!{formula}s formed from them using propositional
connectives.  We say that propositional logic is \emph{truth
  functional}, because its semantics is given by functions of truth
values. In particular, in propositional logic we leave out of
consideration any further determination of truth and falsity, e.g.,
whether something is necessarily true rather than just contingently
true, or whether something is known to be true, or whether something
is true now rather than was true or will be true.  We only consider
two truth values true ($\True$) and false ($\False$), and so exclude
from discussion the possibility that a statement may be neither true
nor false, or only half true. We also concentrate only on connectives where
the truth value of !!a{formula} built from them is completely
determined by the truth values of its parts (and not, say, on its
meaning). In particular, whether the truth value of conditionals in
English is truth functional in this sense is contentious. The material
conditional~$\lif$ is; other logics deal with conditionals that are
not truth functional.

In order to develop the theory and metatheory of truth-functional
propositional logic, we must first define the syntax and semantics of
its expressions.  We will describe one way of constructing
!!{formula}s from !!{propositional variable}s using the connectives.
Alternative definitions are possible. Other systems will choose
different symbols, will select different sets of connectives as
primitive, and will use parentheses differently (or even not at all,
as in the case of so-called Polish notation).  What all approaches
have in common, though, is that the formation rules define the set of
!!{formula}s \emph{inductively}. If done properly, every expression
can result essentially in only one way according to the formation
rules.  The inductive definition resulting in expressions that are
\emph{uniquely readable} means we can give meanings to these
expressions using the same method---inductive definition.

Giving the meaning of expressions is the domain of semantics.  The
central concept in semantics for propositional logic is that of
satisfaction in !!a{valuation}. !!^a{valuation}~$\pAssign{v}$ assigns
truth values $\True$, $\False$ to the !!{propositional variable}s. Any
!!{valuation} determines a truth value $\pValue{v}(!A)$ for any
!!{formula}~$!A$.  !!^a{formula} is satisfied in
!!a{valuation}~$\pAssign{v}$ iff $\pValue{v}(!A) = \True$---we write
this as $\pSat{v}{!A}$. This relation can also be defined by induction on
the structure of~$!A$, using the truth functions for the logical
connectives to define, say, satisfaction of $!A \land !B$ in terms of
satisfaction (or not) of $!A$ and~$!B$.

On the basis of the satisfaction relation $\pSat{v}{!A}$ for sentences
we can then define the basic semantic notions of tautology,
entailment, and satisfiability.  !!^a{formula} is a tautology,
$\Entails !A$, if every !!{valuation} satisfies it, i.e.,
$\pValue{v}(!A) = \True$ for any~$\pAssign{v}$. It is entailed by a
set of !!{formula}s, $\Gamma \Entails !A$, if every !!{valuation} that
satisfies all the !!{formula}s in~$\Gamma$ also satisfies~$!A$. And a
set of !!{formula}s is satisfiable if some !!{valuation} satisfies all
!!{formula}s in it at the same time.  Because !!{formula}s are
inductively defined, and satisfaction is in turn defined by induction
on the structure of !!{formula}s, we can use induction to prove
properties of our semantics and to relate the semantic notions
defined.






\olfileid{pl}{syn}{fml}

\olsection{Propositional \usetoken{P}{formula}}

!!^{formula}s of propositional logic are built up from
\emph{!!{propositional
    variable}s} and the
  propositional constant~$\lfalse$ using \emph{logical
  connectives}.

\begin{enumerate}
\item !!^a{denumerable} set~$\PVar$ of !!{propositional variable}s $\Obj p_0$,
  $\Obj p_1$, \dots
The propositional constant for !!{falsity}~$\lfalse$.

\item The logical connectives:
  \startycommalist
  \ycomma $\lnot$ (negation)
  \ycomma $\land$ (conjunction)
  \ycomma $\lor$ (disjunction)
  \ycomma $\lif$ (!!{conditional})
  
\item Punctuation marks: (, ), and the comma.
\end{enumerate}

We denote this language of propositional logic by $\Lang L_0$.


In addition to the primitive connectives introduced
above, we also use the following \emph{defined} symbols:
\startycommalist
  
  
  
  
  \ycomma $\liff$ (!!{biconditional})
  
  \ycomma $\ltrue$ (!!{truth}).


\begin{explain}
A defined symbol is not officially part of the language, but is
introduced as an informal abbreviation: it allows us to abbreviate
formulas which would, if we only used primitive symbols, get quite
long.  This is obviously an advantage.  The bigger advantage, however,
is that proofs become shorter.  If a symbol is primitive, it has to be
treated separately in proofs. The more primitive symbols, therefore,
the longer our proofs.
\end{explain}




\begin{intro}
You may be familiar with different terminology and symbols than the
ones we use above. Logic texts (and teachers) commonly use either
$\sim$, $\neg$, and~!\ for ``negation'', $\wedge$, $\cdot$, and $\&$
for ``conjunction''.  Commonly used symbols for the ``conditional'' or
``implication'' are $\rightarrow$, $\Rightarrow$, and $\supset$.
Symbols for ``biconditional,'' ``bi-implication,''
  or ``(material) equivalence'' are $\leftrightarrow$,
  $\Leftrightarrow$, and $\equiv$.
The $\lfalse$ symbol is variously called
  ``falsity,'' ``falsum,'' ``absurdity,'' or ``bottom.''
The $\ltrue$ symbol is variously called
  ``truth,'' ``verum,'' or ``top.''
\end{intro}

\begin{defn}[Formula]
\ollabel{defn:formulas}
The set~$\Frm[L_0]$ of \emph{!!{formula}s} of propositional logic
is defined inductively as follows:
\begin{enumerate}
$\lfalse$ is an atomic !!{formula}.



\item Every !!{propositional variable}~$\Obj p_i$ is an atomic
  !!{formula}.

If $!A$ is !!a{formula}, then $\lnot !A$ is
  !!a{formula}.

If $!A$ and $!B$ are !!{formula}s, then $(!A \land
  !B)$ is !!a{formula}.

If $!A$ and $!B$ are !!{formula}s, then $(!A \lor !B)$
  is !!a{formula}.

If $!A$ and $!B$ are !!{formula}s, then $(!A \lif !B)$
  is !!a{formula}.



Nothing else is !!a{formula}.
\end{enumerate}
\end{defn}

\begin{explain}
The definition of !!{formula}s is an
\emph{inductive definition}.  Essentially, we construct the set of
!!{formula}s in infinitely many stages.  In the initial stage, we
pronounce all atomic formulas to be formulas; this corresponds to the
first few cases of the definition, i.e., the cases for

$\lfalse$, 
$\Obj p_i$.  ``Atomic !!{formula}''
thus means any !!{formula} of this form.

The other cases of the definition give rules for constructing new
!!{formula}s out of !!{formula}s already constructed.  At the second
stage, we can use them to construct !!{formula}s out of atomic
!!{formula}s.  At the third stage, we construct new formulas from the
atomic formulas and those obtained in the second stage, and so on.  A
!!{formula} is anything that is eventually constructed at such a
stage, and nothing else.
\end{explain}

When writing a formula $(!B \ast !C)$ constructed from $!B$, $!C$
using a two-place connective~$\ast$, we will often leave out the
outermost pair of parentheses and write simply~$!B \ast !C$.


\begin{defn}
Formulas constructed using the defined operators are to be understood
as follows:

\begin{tagenumerate}{defTrue,defFalse,defNot,defOr,defAnd,defIf,defIff,defEx,defAll}

$\ltrue$ abbreviates
  $\lnot\lfalse$.











$!A \liff !B$ abbreviates $(!A \lif !B) \land (!B
  \lif !A)$.
\end{tagenumerate}
\end{defn}


\begin{defn}[Syntactic identity]
The symbol $\ident$ expresses syntactic identity between strings of
symbols, i.e., $!A \ident !B$ iff $!A$ and $!B$ are strings of symbols
of the same length and which contain the same symbol in each place.
\end{defn}

The $\ident$ symbol may be flanked by strings obtained by
concatenation, e.g., $!A \ident (!B \lor !C)$ means: the string of
symbols~$!A$ is the same string as the one obtained by concatenating
an opening parenthesis, the string $!B$, the $\lor$ symbol, the
string~$!C$, and a closing parenthesis, in this order. If this is the
case, then we know that the first symbol of $!A$ is an opening
parenthesis, $!A$ contains $!B$ as a substring (starting at the second
symbol), that substring is followed by $\lor$, etc.





\olfileid{pl}{syn}{pre}

\olsection{Preliminaries}

\begin{thm}[\emph{Principle of induction on !!{formula}s}]
  \ollabel{thm:induction}  
  If some property~$P$ holds for all the atomic !!{formula}s and is
  such that
  \begin{enumerate}
    it holds for $\lnot !A$ whenever it holds
      for~$!A$;
    it holds for $(!A \land !B)$
      whenever it holds for $!A$ and~$!B$;
    it holds for $(!A \lor !B)$
      whenever it holds for $!A$ and~$!B$;
    it holds for $(!A \lif !B)$
      whenever it holds for $!A$ and~$!B$;
    
  \end{enumerate}
  then $P$ holds for all !!{formula}s.
\end{thm}

\begin{proof}
  Let $S$ be the collection of all !!{formula}s with
  property~$P$. Clearly $S \subseteq \Frm[L_0]$. $S$~satisfies all the
  conditions of \olref[fml]{defn:formulas}: it contains all atomic
  !!{formula}s and is closed under the !!{operator}s.  $\Frm[L_0]$ is
  the smallest such class, so $\Frm[L_0] \subseteq S$. So $\Frm[L_0] = S$, and
  every formula has property~$P$.
\end{proof}

\begin{prop}\ollabel{prop:balanced}
   Any !!{formula} in~$\Frm[L_0]$ is \emph{balanced}, in that it has
   as many left parentheses as right ones.
\end{prop}

\begin{prob}
Prove \olref[pl][syn][pre]{prop:balanced}
\end{prob}

\begin{prop} \ollabel{prop:noinit}
No proper initial segment of !!a{formula} is !!a{formula}.
\end{prop}

\begin{prob}
  Prove \olref[pl][syn][pre]{prop:noinit}
\end{prob}

\begin{prop}[Unique Readability]
Any !!{formula}~$!A$ in $ \Frm[L_0]$ has exactly one parsing as one of
the following
\begin{enumerate}
$\lfalse$.



\item $\Obj p_n$ for some $\Obj p_n \in  \PVar$.
  
$\lnot !B$ for some !!{formula}~$!B$.

$(!B \land !C)$ for some !!{formula}s $!B$ and~$!C$.

$(!B \lor !C)$ for some !!{formula}s $!B$ and~$!C$.

$(!B \lif !C)$ for some !!{formula}s $!B$ and~$!C$.


\end{enumerate}
Moreover, this parsing is \emph{unique}.
\end{prop}

\begin{proof}
By induction on $!A$. For instance, suppose that $!A$ has two distinct
readings as $(!B \lif !C)$ and $(!B' \lif !C')$. Then $!B$ and $!B'$
must be the same (or else one would be a proper initial segment of the
other); so if the two readings of $!A$ are distinct it must be because
$!C$ and $!C'$ are distinct readings of the same sequence of symbols,
which is impossible by the inductive hypothesis.
\end{proof}


\begin{defn}[Uniform Substitution]
If $!A$ and $!B$ are !!{formula}s, and $\Obj p_i$ is a !!{propositional
variable}, then $\Subst{!A}{!B}{\Obj p_i}$ denotes the result of
replacing each occurrence of $\Obj p_i$ by an occurrence of $!B$ in $!A$;
similarly, the simultaneous substitution of $\Obj p_1$, \dots,~$\Obj p_n$ by
!!{formula}s $!B_1$, \dots,~$!B_n$ is denoted by
$\SSubst{!A}{\subst{!B_1}{\Obj p_1},\dots,\subst{!B_n}{\Obj p_n}}$.
\end{defn}

\begin{prob} For each of the five !!{formula}s below determine whether the 
 !!{formula} can be expressed as a substitution \( \Subst{!A}{!B}{\Obj p_i} \)
 where \( !A \) is (i) \( \Obj p_0 \); (ii) \( ( \lnot \Obj p_0 \land \Obj
 p_1) \); and (iii) \( ( ( \lnot \Obj p_0 \lif \Obj p_1 ) \land \Obj
 p_2 ) \). In each case specify the relevant substitution.
  \begin{enumerate}
    \item \( \Obj p_1 \) 
    \item \( ( \lnot \Obj p_0 \land \Obj p_0 ) \)
    \item \( ( ( \Obj p_0 \lor \Obj p_1 ) \land \Obj p_2 )  \)
    \item \( \lnot ( ( \Obj p_0 \lif \Obj p_1 ) \land \Obj p_2 ) \)
    \item \( (( \lnot ( \Obj p_0 \lif \Obj p_1 ) \lif ( \Obj p_0 \lor \Obj p_1 )) \land \lnot ( \Obj p_0 \land \Obj p_1 )) \)
  \end{enumerate}
\end{prob}

\begin{prob}
Give a mathematically rigorous definition of $\Subst{!A}{!B}{p}$ by
induction.
\end{prob}





\olfileid{pl}{syn}{fseq}

\olsection{Formation Sequences}

Defining !!{formula}s via an inductive definition, and the
complementary technique of proving properties of !!{formula}s via
induction, is an elegant and efficient approach. However, it can
also be useful to consider a more bottom-up, step-by-step approach
to the construction of !!{formula}s, which we do here using the
notion of a \emph{formation sequence}.

\begin{defn}[Formation sequences for formulas]
\ollabel{defn:fseq-frm}
A finite sequence $\tuple{!A_0,\dotsc,!A_n}$ of strings of
symbols from the language~$\Lang L_0$ is a \emph{formation
sequence} for $!A$ if $!A \ident !A_n$ and for all $i \leq n$,
either $!A_i$ is an atomic formula or there exist $j,k < i$
such that one of the following holds:
\begin{enumerate}
    $!A_i \ident \lnot !A_j$.
    $!A_i \ident (!A_j \land !A_k)$.
    $!A_i \ident (!A_j \lor !A_k)$.
    $!A_i \ident (!A_j \lif !A_k)$.
    
\end{enumerate}
\end{defn}

\begin{ex}
\[
\tuple{
    \Obj p_0,
    \Obj p_1,
    (\Obj p_1 \land \Obj p_0),
    \lnot (\Obj p_1 \land \Obj p_0)
}
\]
is a formation sequence of
$\lnot (\Obj p_1 \land \Obj p_0)$, as is
\[
\tuple{
    \Obj p_0,
    \Obj p_1,
    \Obj p_0,
    (\Obj p_1 \land \Obj p_0),
    (\Obj p_0 \lif \Obj p_1),
    \lnot (\Obj p_1 \land \Obj p_0)
}.
\]

As can be seen from the second example, formation sequences
may contain `junk': formulas which are redundant or do not
contribute to the construction.
\end{ex}

\begin{prop}
\ollabel{prop:formed}
Every !!{formula}~$!A$ in~$\Frm[L_0]$ has a formation sequence.
\end{prop}

\begin{proof}
Suppose $!A$ is atomic. Then the sequence~$\tuple{!A}$ is a
formation sequence for~$!A$.

Now suppose that $!B$ and~$!C$ have formation sequences
$\tuple{!B_0,\dotsc,!B_n}$ and $\tuple{!C_0,\dotsc,!C_m}$
respectively.

\begin{enumerate}
  If $!A \ident \lnot !B$,
      then $\tuple{!B_0,\dotsc,!B_n,\lnot !B_n}$
      is a formation sequence for~$!A$.
  If $!A \ident (!B \land !C)$,
      then $\tuple{!B_0,\dotsc,!B_n,!C_0,\dotsc,!C_m,(!B_n \land !C_m)}$
      is a formation sequence for~$!A$.
  If $!A \ident (!B \lor !C)$,
      then $\tuple{!B_0,\dotsc,!B_n,!C_0,\dotsc,!C_m,(!B_n \lor !C_m)}$
      is a formation sequence for~$!A$.
  If $!A \ident (!B \lif !C)$,
      then $\tuple{!B_0,\dotsc,!B_n,!C_0,\dotsc,!C_m,(!B_n \lif !C_m)}$
      is a formation sequence for~$!A$.
  
\end{enumerate}
By the principle of induction on !!{formula}s,
every !!{formula} has a formation sequence.
\end{proof}

We can also prove the converse. This is important because it shows
that our two ways of defining formulas are equivalent: they give
the same results. It also means that we can prove theorems about
formulas by using ordinary induction on the length of formation
sequences.

\begin{lem}
\ollabel{lem:fseq-init}
Suppose that $\tuple{!A_0,\dotsc,!A_n}$ is a formation sequence
for~$!A_n$, and that $k \leq n$. Then $\tuple{!A_0,\dotsc,!A_k}$
is a formation sequence for~$!A_k$.
\end{lem}

\begin{proof}
Exercise.
\end{proof}

\begin{thm}
\ollabel{thm:fseq-frm-equiv}
$\Frm[L_0]$ is the set of all strings of symbols
in the language~$\Lang L_0$ with a formation sequence.
\end{thm}

\begin{proof}
Let $F$ be the set of all strings of symbols in the
language~$\Lang L_0$ that have a formation sequence.
We have seen in \olref[pl][syn][fseq]{prop:formed} that
$\Frm[L_0] \subseteq F$, so now we prove the converse.

Suppose $!A$ has a formation sequence $\tuple{!A_0,\dotsc,!A_n}$.
We prove that $!A \in \Frm[L_0]$ by strong induction on~$n$.
Our induction hypothesis is that every string of symbols with a
formation sequence of length $m < n$ is in~$\Frm[L_0]$.
By the definition of a formation sequence, either $!A_n$ is
atomic or there must exist $j,k < n$ such that one of the
following is the case:
\begin{enumerate}
    $!A_n \ident \lnot !A_j$.
    $!A_n \ident (!A_j \land !A_k)$.
    $!A_n \ident (!A_j \lor !A_k)$.
    $!A_n \ident (!A_j \lif !A_k)$.
    
\end{enumerate}
Now we reason by cases. If $!A_n$ is atomic then
$!A_n \in \Frm[L_0]$. Suppose instead that $!A \equiv
(!A_j \land !A_k)$. By \olref[pl][syn][fseq]{lem:fseq-init},
$\tuple{!A_0,\dotsc,!A_j}$ and $\tuple{!A_0,\dotsc,!A_k}$ are
formation sequences for $!A_j$ and~$!A_k$ respectively. Since
these are proper initial subsequences of the formation sequence
for~$!A$, they both have length less than~$n$. Therefore by
the induction hypothesis, $!A_j$ and~$!A_k$ are in~$\Frm[L_0]$,
and so by the definition of !!a{formula}, so is
$(!A_j \land !A_k)$. The other cases follow by parallel
reasoning.
\end{proof}





\olfileid{pl}{syn}{val}

\olsection{\usetoken{P}{valuation} and Satisfaction}

\begin{defn}[!!^{valuation}s]
Let $\{\True, \False\}$ be the set of the two truth values, ``true''
and ``false.'' A \emph{!!{valuation}} for $\Lang{L_0}$ is a
function~$\pAssign{v}$ assigning either $\True$ or $\False$ to the
!!{propositional variable}s of the language, i.e., $\pAssign{v} \colon
\PVar \to \{\True, \False \}$.
\end{defn}

\begin{defn}
  Given !!a{valuation}~$\pAssign{v}$, define the evaluation function
  $\pValue{v} \colon \Frm[L_0] \to \{\True, \False \}$ inductively by:
  \begin{align*}
    \pValue{v}(\lfalse) & = \False; \\
    
    \pValue{v}(\Obj p_n) & = \pAssign{v}(\Obj p_n); \\
    \pValue{v}(\lnot !A) & = \begin{cases}
        \True & \text{if } \pValue{v}(!A) = \False;\\
        \False & \text{otherwise.}
      \end{cases} \\ 
    \pValue{v}(!A \land !B) & = \begin{cases}
        \True &
        \text{if $\pValue{v}(!A) = \True$ and $\pValue{v}(!B) = \True$;}\\
        \False &
        \text{if $\pValue{v}(!A) = \False$ or $\pValue{v}(!B) = \False$}.
      \end{cases}\\
    \pValue{v}(!A \lor !B) & = \begin{cases}
        \True &
        \text{if $\pValue{v}(!A) = \True$ or $\pValue{v}(!B) = \True$;}\\
        \False &
        \text{if $\pValue{v}(!A) = \False$ and $\pValue{v}(!B) = \False$}.
      \end{cases}\\
    \pValue{v}(!A \lif !B) & = \begin{cases}
        \True &
        \text{if $\pValue{v}(!A) = \False$ or $\pValue{v}(!B) = \True$;}\\
        \False &
        \text{if $\pValue{v}(!A) = \True$ and $\pValue{v}(!B) = \False$}.
      \end{cases}\\
    
\end{align*}
\end{defn}

\begin{explain}
The clauses correspond to the following truth tables:
\begin{center}

  \begin{tabular}{|c||c|} \hline
    $!A$ & $ \lnot !A$ \\
    \hline \hline
    $\True$ & $\False$ \\
    $\False$ & $\True$ \\
    \hline
  \end{tabular}


  \begin{tabular}{|cc||c|} \hline
    $!A$ & $!B$ & $!A \land !B$ \\
    \hline \hline
    $\True$ & $\True$ & $\True$ \\
    $\True$ & $\False$ & $\False$ \\
    $\False$ & $\True$ & $\False$ \\
    $\False$ & $\False$ & $\False$ \\
    \hline
  \end{tabular}


  \begin{tabular}{|cc||c|} \hline
    $!A$ & $!B$ & $!A \lor !B$ \\
    \hline \hline
    $\True$ & $\True$ & $\True$ \\
    $\True$ & $\False$ & $\True$ \\
    $\False$ & $\True$ & $\True$ \\
    $\False$ & $\False$ & $\False$ \\
    \hline
  \end{tabular}

\\[1em]

  \begin{tabular}{|cc||c|} \hline
    $!A$ & $!B$ & $!A \lif !B$ \\
    \hline \hline
    $\True$ & $\True$ & $\True$ \\
    $\True$ & $\False$ & $\False$ \\
    $\False$ & $\True$ & $\True$ \\
    $\False$ & $\False$ & $\True$ \\
    \hline
  \end{tabular}


\end{center}
\end{explain}

\begin{prob}
Consider adding to $\Lang{L_0}$ a ternary connective $\diamondsuit$
with evaluation given by
\begin{gather*}
  \pValue{v}(\diamondsuit( !A , !B , !C ) ) = \begin{cases}
    \pValue{v}( !B ) &
    \text{if $\pValue{v}(!A) = \True$;}\\
    \pValue{v}( !C ) &
    \text{if $\pValue{v}(!A) = \False $}.
  \end{cases}
\end{gather*}
Write down the truth table for this connective.
\end{prob}

\begin{thm}[Local Determination]
\ollabel{thm:LocalDetermination} Suppose that $\pAssign{v_1}$ and
$\pAssign{v_2}$ are !!{valuation}s that agree on the !!{propositional
variable}s occurring in $!A$, i.e., $\pAssign{v_1}(\Obj p_n) =
\pAssign{v_2}(\Obj p_n)$ whenever $\Obj p_n$ occurs in some
!!{formula}~$!A$. Then $\pValue{v_1}$ and $\pValue{v_2}$ also agree
on~$!A$, i.e., $\pValue{v_1}(!A) = \pValue{v_2}(!A)$.
\end{thm}

\begin{proof}
By induction on $!A$.
\end{proof}

\begin{defn}[Satisfaction]
\ollabel{defn:satisfaction} We can inductively define the notion of
  \emph{satisfaction of !!a{formula}~$!A$ by
  !!a{valuation}~$\pAssign{v}$}, $\pSat{v}{!A}$, as follows.
  (We write $\pSat/{v}{!A}$ to mean ``not $\pSat{v}{!A}$.'')
\begin{enumerate}

  \indcase{!A}{\lfalse}{$\pSat/{v}{\indfrm}$.}



\item \indcase{!A}{\Obj p_i}{$\pSat{v}{\indfrm}$
  iff $\pAssign{v}(\Obj p_i) = \True$.}


  \indcase{!A}{\lnot !B}{$\pSat{v}{\indfrm}$ iff
    $\pSat/{v}{!B}$.}


  \indcase{!A}{(!B \land !C)}{$\pSat{v}{\indfrm}$ iff $\pSat{v}{!B}$
    and $\pSat{v}{!C}$.}


  \indcase{!A}{(!B \lor !C)}{$\pSat{v}{\indfrm}$ iff
    $\pSat{v}{!B}$ or $\pSat{v}{!C}$ (or both).}


  \indcase{!A}{(!B \lif !C)}{$\pSat{v}{\indfrm}$ iff $\pSat/{v}{!B}$
    or $\pSat{v}{!C}$ (or both).}


\end{enumerate}
If $\Gamma$ is a set of !!{formula}s, $\pSat{v}{\Gamma}$ iff
$\pSat{v}{!A}$ for every~$!A \in \Gamma$.
\end{defn}

\begin{prop}\ollabel{prop:sat-value}
  $\pSat{v}{!A}$ iff $\pValue{v}(!A) = \True$.
\end{prop}

\begin{proof}
  By induction on~$!A$.
\end{proof}

\begin{prob}
  Prove \olref[pl][syn][val]{prop:sat-value}
\end{prob}





\olfileid{pl}{syn}{sem}

\olsection{Semantic Notions}

We define the following semantic notions:

\begin{defn} 
\begin{enumerate}
\item !!^a{formula}~$!A$ is \emph{satisfiable} if for
  some~$\pAssign{v}$, $\pSat{v}{!A}$; it is
  \emph{unsatisfiable} if for no $\pAssign{v}$, $\pSat{v}{!A}$;
\item !!^a{formula}~$!A$ is a \emph{tautology} if $\pSat{v}{!A}$ for
  all !!{valuation}s~$\pAssign{v}$;
\item !!^a{formula}~$!A$ is \emph{contingent} if it is satisfiable but
  not a tautology;
\item If $\Gamma$ is a set of !!{formula}s, $\Gamma \Entails !A$ (``$\Gamma$
  entails $!A$'') if and only if $\pSat{v}{!A}$ for every
  !!{valuation}~$\pAssign{v}$ for which $\pSat{v}{\Gamma}$.
\item If $\Gamma$ is a set of !!{formula}s, $\Gamma$ is
  \emph{satisfiable} if there is !!a{valuation}~$\pAssign{v}$ for which
  $\pSat{v}{\Gamma}$, and $\Gamma$ is
  \emph{unsatisfiable} otherwise.
\end{enumerate} 
\end{defn}

\begin{prob} 
 For each of the following four !!{formula}s determine whether it
 is (a)~satisfiable, (b)~tautology, and (c)~contingent.
\begin{enumerate}
  \item \( ( \Obj p_0 \lif ( \lnot \Obj p_1 \lif \lnot \Obj p_0 ) ) \).
  \item \( ( ( \Obj p_0 \land \lnot \Obj p_1 ) \lif ( \lnot \Obj p_0 \land \Obj p_2 )) \liff ( ( \Obj p_2 \lif \Obj p_0 ) \lif ( \Obj p_0 \lif \Obj p_1 )) \).
  \item \( ( \Obj p_0 \liff \Obj p_1 ) \lif ( \Obj p_2 \liff \lnot \Obj p_1 ) \).
  \item \( (( \Obj p_0 \liff ( \lnot \Obj p_1 \land \Obj p_2 )) \lor ( \Obj p_2 \lif ( \Obj p_0 \liff \Obj p_1 ))) \).
\end{enumerate}
\end{prob}

\begin{prop}
\ollabel{prop:semanticalfacts} 
\begin{enumerate} 
\item $!A$ is a tautology if and only if
  $\emptyset \Entails !A$; 
\item If $\Gamma \Entails !A$ and $\Gamma \Entails !A \lif !B$ then
  $\Gamma \Entails !B$;
\item If $\Gamma$ is satisfiable then every finite subset of $\Gamma$
  is also satisfiable; 
\item \ollabel{def:monotonicity} Monotonicity: if $\Gamma \subseteq \Delta$
  and $\Gamma \Entails !A$ then also $\Delta \Entails !A$;
\item \ollabel{def:Cut} Transitivity: if $\Gamma \Entails !A$ and
  $\Delta \cup \{ !A\} \Entails !B$ then $\Gamma \cup \Delta \Entails
  !B$.
\end{enumerate}
\end{prop}

\begin{proof}
Exercise.
\end{proof}

\begin{prob}
Prove \olref[pl][syn][sem]{prop:semanticalfacts}
\end{prob}

\begin{prop}\ollabel{prop:entails-unsat}
  $\Gamma \Entails !A$ if and only if $\Gamma \cup \{\lnot !A\}$
  is unsatisfiable.
\end{prop}

\begin{proof}
Exercise.
\end{proof}

\begin{prob}
Prove \olref[pl][syn][sem]{prop:entails-unsat}
\end{prob}

\begin{thm}[Semantic Deduction Theorem]
  \ollabel{thm:sem-deduction} $\Gamma \Entails !A \lif !B$ if and only
  if $\Gamma \cup \{!A\} \Entails !B$.
\end{thm}

\begin{proof}
Exercise.
\end{proof}

\begin{prob}
Prove \olref[pl][syn][sem]{thm:sem-deduction}
\end{prob}



\OLEndChapterHook







